\documentclass[11pt]{article}
\usepackage{mathblatt}
\begin{document}
\weit
\typeout{PROBE-BEGIN brueche-dezimalzahlen-e1-k1-s6-v17}
\begin{aufgabe}{\detokenize{brueche-dezimalzahlen-e1-k1-s6-v17}}
\begin{teile}
\teil Probe

P \streifen[0]{20}{}{} \quad Q \kreissektor{60}{} \quad R \bruchrechteck{1}{4} \quad S \bruchkreis{1}{3}
\end{teile}
\end{aufgabe}
\typeout{PROBE-END brueche-dezimalzahlen-e1-k1-s6-v17}
\typeout{PROBE-BEGIN prozentrechnung-e3-k1-s1-v1}
\begin{aufgabe}{\detokenize{prozentrechnung-e3-k1-s1-v1}}
\begin{teile}
\teil Probe

\streifen[2]{50}{$0$}{$80\,€$}
\end{teile}
\end{aufgabe}
\typeout{PROBE-END prozentrechnung-e3-k1-s1-v1}
\typeout{PROBE-BEGIN brueche-dezimalzahlen-e1-k1-s6-v18}
\begin{aufgabe}{\detokenize{brueche-dezimalzahlen-e1-k1-s6-v18}}
\begin{teile}
\teil Probe

P \bruchkreis{1}{4} \quad Q \streifen[0]{30}{}{} \quad R \bruchrechteck{1}{6} \quad S \kreissektor{72}{}
\end{teile}
\end{aufgabe}
\typeout{PROBE-END brueche-dezimalzahlen-e1-k1-s6-v18}
\typeout{PROBE-BEGIN prozentrechnung-e3-k1-s1-v2}
\begin{aufgabe}{\detokenize{prozentrechnung-e3-k1-s1-v2}}
\begin{teile}
\teil Probe

\streifen[4]{25}{$0$}{$60$ kg}
\end{teile}
\end{aufgabe}
\typeout{PROBE-END prozentrechnung-e3-k1-s1-v2}
\typeout{PROBE-BEGIN prozentrechnung-e3-k1-s1-v3}
\begin{aufgabe}{\detokenize{prozentrechnung-e3-k1-s1-v3}}
\begin{teile}
\teil Probe

\streifen{10}{$0$}{$300$ m}
\end{teile}
\end{aufgabe}
\typeout{PROBE-END prozentrechnung-e3-k1-s1-v3}
\end{document}
